\documentclass[11pt]{article}
\usepackage[utf8]{inputenc}
\usepackage{amsmath}
\usepackage{geometry}
\usepackage{listings}
\usepackage{xcolor}

\geometry{letterpaper, margin=1in}

\definecolor{codegreen}{rgb}{0,0.6,0}
\definecolor{backcolour}{rgb}{0.95,0.95,0.92}

\lstdefinestyle{mystyle}{
    backgroundcolor=\color{backcolour},   
    commentstyle=\color{codegreen},
    basicstyle=\ttfamily\footnotesize,
    breaklines=true,                 
    numbers=left,                    
    numbersep=5pt,                  
    tabsize=2
}

\lstset{style=mystyle}

\title{\textbf{Editorial: Problem A - Touko's Perfection}}
\author{Setter: Hudson}
\date{}

\begin{document}

\maketitle

\section*{Problem Overview}
We are given an array $A$. We want to construct a non-decreasing array $B$ such that $\sum |A_i - B_i|$ is minimized.

\section*{Difficulty Analysis}
Rating: \textbf{2100}.
A naive DP approach would define $dp[i][val]$ as the min cost to set the $i$-th element to $val$. This is $O(N \cdot \max(A))$, which is too slow.
This is a standard application of the \textbf{Slope Trick}.

\section*{Solution Approach}
Let $F_i(x)$ be the minimum cost to fix the prefix $1 \dots i$ such that the $i$-th element is exactly $x$.
$$ F_i(x) = |x - A_i| + \min_{y \le x} F_{i-1}(y) $$
The function $F_i(x)$ is always convex. We can maintain this function using a priority queue that stores the locations where the slope of the function changes.
For the condition $B_i \ge B_{i-1}$, the transition $\min_{y \le x} F_{i-1}(y)$ corresponds to keeping the "left" part of the convex function and flattening the "right" part (since we can always choose a larger value for free).
Specifically, if we maintain the change points of the slope in a Max-Priority Queue:
1. Push $A_i$ onto the queue twice.
2. If the maximum element in the queue (let's say $top$) is greater than $A_i$, it means the optimal "valley" of the function is to the left of the current $top$.
3. We pop the top element. The difference $top - A_i$ adds to our minimum cost.

\section*{Implementation (C++)}

\begin{lstlisting}[language=C++]
#include <iostream>
#include <vector>
#include <queue>
#include <algorithm>

using namespace std;

int main() {
    ios_base::sync_with_stdio(false);
    cin.tie(NULL);

    int n;
    if (!(cin >> n)) return 0;

    // Max Priority Queue stores the slope change points
    // Specifically, for the "left" side of the convex function
    priority_queue<long long> pq;
    long long ans = 0;

    for (int i = 0; i < n; ++i) {
        long long a;
        cin >> a;

        // Add the function |x - a|
        // This adds a slope change of +1 at x=a and -1 at x=a
        // In the context of the running min-prefix DP, we push 'a' twice.
        pq.push(a);
        pq.push(a);

        // The transition min_{y <= x} F(y) means we effectively
        // flatten the slope for the right side.
        // We remove the rightmost slope change point if it's too high.
        long long max_el = pq.top();
        pq.pop();

        // If the current optimal point (max_el) is greater than a,
        // it contributes to the cost.
        ans += (max_el - a);
    }

    cout << ans << endl;

    return 0;
}
\end{lstlisting}

\end{document}
